\section{Discussion}\label{sec:discussion}

\textbf{The picture.} The village behaves as a weakly coupled, subcritical paramagnet under strong fields. Each agent is a spin that updates on its own call clock. External fields set most of its state: the scheduler, the goal, the room and the personal style of the agent. Agents couple only through what they read, at the next model call, with a loop gain near 0.13. A read moves the next statements of the reader for a few calls. The content of the agent then relaxes back over about one hour. A new goal erases the previous content state within one day. Rooms split in content without an external cause and lose the split at the next goal. We found no collective variable beyond pairs and fields, except one slow content mode in the early single-room village. Table~\ref{tab:theses} lists the theses with the highest estimated usefulness.

\begin{table*}[t]
\caption{The 20 research theses with the highest estimated usefulness among those with credence $\geq 0.6$ (Claude-judged, Sec.~\ref{sec:method}). $p$: credence (faithfulness); $V$: value if true (0--5); EU $=pV$; M: mechanism depth; $^\ast$fragile. No thesis has passed a confirmation test on reserved data.}
\label{tab:theses}
\scriptsize\renewcommand{\arraystretch}{1.05}
\begin{tabular}{@{}l p{0.70\textwidth} c c c c@{}}
\toprule
ID & thesis (scored claim, short form) & $p$ & $V$ & EU & M\\ \midrule
H44 & A forced context wipe causes orderly re-reading and a 26\% write dip (4-11\% of output) and breaks command loops; it adds no room susceptibility. & 0.88 & 3.7 & 3.23 & M2\\
H50 & Activity co-movement is the scheduler's field; talk spreads by a coupling acting at one read-out call (45/71 units, onset hop 1, none negative). & 0.85 & 3.5 & 2.98 & M1\\
H38 & Joint silences are rare and operator-scheduled, not platform stalls; in the always-on regime 70-80\% of co-activation is the runner's daily start and stop. & 0.90$^\ast$ & 3.2 & 2.93 & M2\\
H08 & Responses start at the call that receives the message (addressing 14/17, replies 17/17, nudges too); a forced erasure cuts reply coupling by 21\%. & 0.91 & 3.2 & 2.88 & M2\\
H15 & Memory loss costs nothing visible in a day; a forced context wipe costs 39\% of committed work for 10 calls, whatever the agent saves. & 0.78$^\ast$ & 3.5 & 2.73 & M1\\
H54 & Day-1 content lands on its kickoff (top-1 18/33), \#51 agents on their private goals (0.95); frozen projects carry the goal's name; specificity sets no spread. & 0.92 & 2.8 & 2.60 & M2\\
H40 & In the computer-use scaffold agents couple per model call, not per minute; per-hour coupling is a family trait times call rate. Regime I breaks this. & 0.72 & 3.5 & 2.52 & M1\\
H46 & Style is not conserved: goal switches (T 0.69), forced erasures (0.56, content flat) and assigned registers move it; it identifies agents at 0.76 (content 0.51). & 0.75 & 3.2 & 2.44 & M1\\
H02 & Activity-timing Ising couplings stay at the null floor on corrected data; the regime-III collective coupling is the operator's day edges. & 0.94$^\ast$ & 2.5 & 2.35 & M1\\
H69 & Restatement loops are copies from the context window (source in context 3.4x, erasure exit 2.4x) but are not triggered by a self-share threshold or ended by novel input. & 0.85 & 2.8 & 2.34 & M1\\
H05 & Rooms decouple chat only through what agents read; the effect is stronger on fixed data, and work is not room-structured. & 0.88 & 2.5 & 2.20 & M1\\
H36 & Fluctuation statistics barely flag swarm reorganizations: a quarter of goal changes, all via content, at the random-date boundary; a topic-shift detector does far better. & 0.80$^\ast$ & 2.8 & 2.20 & M0\\
H67 & Read-out talk coupling is subcritical (g\_lag $\sim$0.13 in regime III, $\sim$0 elsewhere), carried by named messages; the equal-time dial overstates it (lagged > equal-time in only 7/35 periods). & 0.79 & 2.8 & 2.17 & M1\\
H13 & Model families share a stable content field that is writing style; behavior carries a weak family trace; no family homophily; rooms carry coupling. & 0.87 & 2.5 & 2.17 & M1\\
H16 & Traps age by themselves, not by Kramers escape; one directed kick at a pause point helps; real failure loops age only in \#51. & 0.70 & 3.0 & 2.10 & M1\\
H52 & Matched on salience, humans get about twice the reply rate and a small content pull; agent leaders share the premium, and off-role messages get none. & 0.60 & 3.5 & 2.10 & M1\\
H11 & Agents herd onto shared projects whatever the goal, in commits as well as attention; spread appears only as fixed ownership. & 0.75 & 2.8 & 2.06 & M1\\
H29 & Influence is address-dependent (sharper on the context ledger); the content-pull driver ranking fails, a reply-graph ranking weakly predicts spread; operator messages do not relay. & 0.65 & 3.2 & 2.06 & M1\\
H74 & A daily schema diff of log records dates platform and provider API changes to the day with zero placebo false alarms; the fused five-channel alarm is mixed (AUC 0.68) and too noisy (FAR 0.18). & 0.63 & 3.2 & 2.05 & M1\\
H28 & Links co-move with project switches (HR 2.6), but recipients switch before they can read them: links mark bursts; removal trims peaks at most 30\%. & 0.81 & 2.5 & 2.03 & M1\\
\bottomrule
\end{tabular}
\end{table*}


\textbf{What transfers.} Three tools need only standard logs. The partition contrast (read against in flight, at the same lag) separates coupling from common drives. A swarm can use it if it logs what each agent saw. The Fano sum rule $\Phi=1/(1-g)^2$ checks the read-out gain from talk counts, with no free parameter. An excess shows a field or a coupling that the read-out loop does not include. The impostor checklist and the synthetic runs found estimator errors before real data in most hypotheses. Several apparent effects disappeared under these checks: a nudge Onsager violation, an attention exponent and a collective entropy production.

\textbf{Limits.} The results come from one village with one operator, and its scaffold changed three times. No claim has passed a held-out confirmation test. The confirmation scripts are frozen and wait for sign-off. The credences are judgements by Claude under a written rubric, with a blind second judge. They are not frequencies. Claude agents also ran the analyses, so one model family designed, ran and scored the tests. Dated predictions, synthetic checks and frozen scripts reduce this dependence, but they do not remove it. Content measures use two embedding models and one zero-shot labeller, which we checked against blind references. Goal periods longer than about one week drift (H120). A fit over a longer window gives the mean of a changing state. We tested 132 hypotheses, so some passes are chance. We report every hypothesis, including the failures.

\textbf{Next steps.} First, run the confirmation scripts on the reserved periods for the core claims: read-out gating (H08), the call clock (H40), the subcritical gain (H67) and the sum rule (H111). Second, test the open slow mode (H81) again, with the record divided at the dates where the mode fades.

\textbf{Data and code.} We used the AI Village data under the research terms of AI Digest and do not distribute them. The analysis code, cards, scores and figures are in the project repository.
